\section{The eight uniform family parameter bounds}

% BEGIN A78P1
\paragraph{A78P1. The host class number at a family member.}
Let $\mathcal F$ be any original $\mathrm{FamP}(\epsilon)$.
Throughout assume $\mathcal F.x_1<x$, take $i\in\mathcal F.\mathrm{Ix}(x)$,
and write $D,L,P$ for the member, its size datum and its specified
parameter. Then $L.\mathrm{sc}.x=x$, $L.\mathrm{sc}.\epsilon=\epsilon$,
$x>1$, and $0<\epsilon<1/52$. Write
\[
 u_N=2+8\epsilon+\epsilon/100,\qquad u_q=\epsilon/5+\epsilon/100,
 \qquad A=|u_N|+|u_q|+1,\quad B=x^A .
\]
The upper family bounds give $n\le x^{u_N}$ and
$\log q,\log d\le x^{u_q}$. Thus $B\ge1$ dominates all three.
The host field $E$ is CM of degree $2f$, where $f=2n$. Writing
$\delta_E=|\Delta_E|^{1/(2f)}$, the discriminant bounds on this
member are $1\le\delta_E\le(Q_S/2)\sqrt{qd}$ and $Q_S\le n^2$.
Hence
\[
   0\le\log\delta_E
     \le 2\log n+(\log q+\log d)/2\le3B,\qquad f\le2B .
\]

Use the ideal count in OB4 up to its discriminant bound. More
explicitly, Minkowski's theorem represents each ideal class by an
integral ideal of norm at most $Y=\sqrt{|\Delta_E|}$. Its constant
$(4/\pi)^f (2f)!/(2f)^{2f}$ is at most one, since
$(2f)!\le((2f+1)/2)^{2f}$ and $(2f+1)/(4f)\le3/4$.
At every prime $p$ the generating factor for ideals has the form
$\prod_{\mathfrak p\mid p}(1-T^{f_{\mathfrak p}})^{-1}$, bounded
coefficientwise by $(1-T)^{-2f}$. The total number up to $Y$ is
at most the number of positive $2f$-tuples of product at most $Y$.
For $s\ge1$ the latter count in $s$ coordinates is at most
$Y(1+\log Y)^{s-1}$: the case $s=1$ is $\lfloor Y\rfloor\le Y$,
and summing the bound for $Y/a$ over $1\le a\le Y$ uses
$\sum_{a\le Y}1/a\le1+\log Y$. Since $Y\ge1$, we may increase
the exponent by one. Taking logarithms and writing
$v=\log Y=f\log\delta_E\le6B^2$ gives
\[
   \log h_E\le v+2f\log(1+v)
   \le (1+2f)v\le30B^3\le C_E x^{C_E},
   \qquad C_E=30+3A .
\]
This last estimate uses $1+2f\le5B$ and $x\ge1$. Its constant
depends only on the family exponent.
% END A78P1

% BEGIN A78P2
\paragraph{A78P2. Counting the finite group for the restricted rows.}
Let $G$ denote the group's positive cardinal
$D.\mathrm{what}_2.\mathrm{restRowsFiniteGroupCard}$. Set
$U=x^\epsilon\ge1$ and $u_K=9/2+10\epsilon+\epsilon/100$.
For $x\ge8$, the bounds $k\le x^{u_K}$, $n\le x^{u_N}$ and
$0<\epsilon<1/52$ give $k\le x^5$ and $f=2n\le2x^{5/2}\le x^3$.
The rank definition gives $r\le2k$. For each of the two auxiliary
primes $g_j$, the family size bound gives
$\log g_j\le x^{\epsilon/100}U\le x$. The level-cardinality
identity expresses $G$ as
\[
     \prod_{j=0,1}|\mathrm{GL}_r(\mathbb Z/g_j\mathbb Z)|^f .
\]
Forgetting the inverse injects this group at a single level into
the $g_j^{\,r^2}$ matrices. As a result
\[
 \log G\le r^2f(\log g_0+\log g_1)
          \le (2x^5)^2x^3(2x)
          =8x^{14}\le x^{15}\le15x^{15}U .
\]
This gives the required constant $C_G=15$. All the inequalities
hold at the same numerical cutoff eight, independently of $i$.
% END A78P2

% BEGIN A78P3
\paragraph{A78P3. A single exponent for all eight entries.}
Write $p_a$ for entry $a$ of the family's fixed list $\mathrm{PK}$
in zero-based numbering, with default zero. The identity
$\operatorname{parConsts}(P)=\mathrm{PK}$ gives $P.C_{L_w}=p_7$
and $P.C_{YL}=p_{19}$. Let $u_t=\mathcal F.C_t+\epsilon/100$ and take
\[
 C_1=\max\{1,p_7,p_{19},u_K,u_N,u_t,C_E,15\}.
\]
The original size conditions, with $L.C_t=\mathcal F.C_t$, give
$k\le x^{u_K}$, $n\le x^{u_N}$, $t\le x^{u_t}$, and
$x^{-\epsilon}\le L.\mathrm{de2}$. For $x\ge8$, A78P1 and A78P2
give the two logarithmic bounds. If $c\le C$ and $C\ge0$, then
$c x^c\le Cx^c\le Cx^C$ because $x\ge1$. For the order of the
finite group we may also multiply this comparison by $U\ge0$.
Thus the eight fields of the original predicate on this $P$ are
\[
 \begin{gathered}
 P.C_{L_w}\le C_1,\qquad P.C_{YL}\le C_1,\qquad
 k\le x^{C_1},\qquad n\le x^{C_1},\qquad t\le x^{C_1},\\
 x^{-\epsilon}\le L.\mathrm{de2},\qquad
 \log h_E\le C_1x^{C_1},\qquad \log G\le C_1x^{C_1}U .
 \end{gathered}
\]
This proves $\mathrm{admissible\_parameters\_family\_const}$ with
$C_1$ and the cutoff chosen before $x,i$ and its scale guard.
The source's common $C_s$ is at least $C_1$. The same monotonicity
increases the two list bounds, the three rank powers and the two
logarithmic bounds to $C_s$. The lower bound on $\mathrm{de2}$ is
independent of the exponent being increased. Including the
parameter cutoff in the final family maximum therefore gives the
original $\mathrm{Concl.par}$ at $C_s$. There is no dependence of
this argument on the hub chosen for the other fields.
% END A78P3
